% Supporting geometry and figure data retained outside the main paper.
% This is a source fragment using the manuscript preamble, notation and
% bibliography. It is not input by main.tex and is not a second document.

\section*{Matched Spherical and Capped Profiles}
\label{app:comparison-profile}

The two walls in Figure~\ref{fig:reduced-patch} use
\begin{equation}
 v_{\rm sph}=W(x^2+y^2+z^2-1),\qquad
 v_{\rm cap}=W(x^2+\psi(y^2+z^2)-1).
 \label{eq:supp-matched-figure-profiles}
\end{equation}
For the cap, take \(L=b=\kappa=1\), \(c=1/2\), \(a=3/5\), and
\[
 \psi(s)=\frac{H((s-a^2)/d)}{H((1-a^2)/d)},\qquad d=1/5,
 \qquad
 H(u)=\begin{cases}\displaystyle\int_0^u e^{-1/z}\,dz&u>0,\\0&u\leq0.
 \end{cases}
\]
Thus \(\psi=0\) for \(\rho\leq a\), \(\psi(1)=1\), and the cap joins
its curved wall smoothly.  The spatial panels show the sections \(z=0\).
Both profiles have exactly the same axial shift \(W(x^2-1)\).

Choose \(W\) smooth and nonincreasing, equal to zero for
\(\phi\leq-0.98\), to \(-1-\phi/2\) on \(|\phi|\leq0.9\), and to
\(-2\) for \(\phi\geq1.8\).  Smooth monotone interpolation gives such a
profile; the Supplement supplies an explicit choice.  Its unique
stationary-limit level is \(\phi=0\).
The front and rear axial roots are \(x=\pm1\), with \(|v_x|=1\).
The numerical traveler in panel (a) starts at the bubble center, at
\(x=\rho=0\), and follows an axial timelike geodesic into the affine cap
collar and through the crossing.  Its interior trajectory and the orange
center worldline use the global profile \(W\), matched to the regular cap
coordinates in the collar.  The plotting script checks normalization,
conserved energy, matching, and the crossing before drawing the figure.
Amber shows only the added outward sector \(p>0,q<0\); the continued
generator \(p=0\) bounds it, and its dotted edges are chart cuts.

\subsection*{Wall shape and matter source}

For a regular front with \(v_x=-\kappa<0\), write \(x=F(y,z)\).
Differentiating its wall equation twice at a stationary generator gives
\begin{equation}
  F_{AB}=\frac{H_{AB}}{\kappa}.
  \label{eq:wall-graph-hessian}
\end{equation}
In principal transverse coordinates for a smooth convex front,
\begin{equation}
  F_{yy}=-C_1,\qquad F_{zz}=-C_2,\qquad C_i>0,
\end{equation}
and the spherical result is \(C_1=C_2=1/R_F\).  An isolated zero of \(H\) is
insufficient: Theorem~\ref{thm:extension-rigidity} requires an open planar cap
with constant \(\kappa\).

If the Einstein equation is imposed, the same wall curvature determines the
null energy component along the increasingly boosted light direction (in
geometrized units):
\begin{equation}
  8\pi T_{KK}=G_{KK}=q^2\operatorname{tr}H,
  \qquad
  G_{KK}=-\kappa(C_1+C_2)q^2<0
  \quad\text{at a convex future front}.
  \label{eq:null-source}
\end{equation}
Thus the transverse wall curvature that produces the tidal obstruction also
gives an unbounded negative null energy component at a convex front.  No energy
condition enters the extension proof: at a saddle point the full tidal tensor
may be nonzero even when its trace, and hence this source component, vanishes.
This local source projection is distinct from the integrated matter
constraints studied in Refs.~\cite{PfenningFord1997,LoboVisser2004}.


\subsection*{Transverse geometry at the aperture}
\label{app:meridional-diagrams}

For the axisymmetric examples, cylindrical coordinates give
\begin{equation}
  g=g_3+\rho^2d\phi^2,\qquad
  g_3=-dt^2+[dx-v(x,\rho)dt]^2+d\rho^2.
  \label{eq:meridional-quotient}
\end{equation}
The spatial panels of Figure~\ref{fig:reduced-patch} show sections of the
wall rather than a conformal diagram of \(g_3\).  Every four-dimensional causal curve projects
to a causal curve of \(g_3\), and every causal curve of \(g_3\) lifts at
constant \(\phi\).  Away from the axis, zero angular momentum geodesics
remain at constant \(\phi\).  The coordinate edge \(\rho=0\) is the symmetry axis, not a
physical boundary.

Use the axial shift \(v_0(x)=v(x,0)\) to define \(p,q\) as above, also away
from the axis.  Put \(\delta=v(x,\rho)-v_0(x)\), and write
\begin{align}
  \theta=dt&=\frac{v_0+1}{2\kappa p}\,dp
                +\frac{v_0-1}{2\kappa q}\,dq,\notag\\
  \eta=dx-v_0dt&=-\frac{v_0+1}{2\kappa p}\,dp
                +\frac{v_0-1}{2\kappa q}\,dq.
  \label{eq:meridional-coframe}
\end{align}
Direct pullback gives the full three-dimensional metric
\begin{equation}
  g_3=C(pq)\,dp\,dq+d\rho^2-2\delta\,\theta\eta
                       +\delta^2\theta^2,
  \qquad \det g_3=-C^2/4.
  \label{eq:meridional-pullback}
\end{equation}
Inside the cap, \(\delta=0\), and its conformal representative in the
axial conformal coordinates is
\begin{equation}
  \Omega^{-2}g_3=-dT^2+dX^2+\Omega^{-2}d\rho^2,
  \qquad \Omega^2=\frac{3C(pq)(1+p^2)(1+q^2)}4.
  \label{eq:meridional-conformal-metric}
\end{equation}
The transverse cone opening depends on \(\Omega\); the geometry is not an
extrusion of a Minkowski diamond.  Outside the cap the additional terms in
Eq.~\eqref{eq:meridional-pullback} must also be retained.  At fixed
\(\rho>a\), with \(\delta_H=v(L,\rho)+1<0\),
\[
  q^2(g_3)_{qq}\longrightarrow
  \frac{\delta_H^2-2\delta_H}{\kappa^2}>0.
\]
For the spherical bubble this happens at every fixed \(\rho>0\); for the
capped bubble it happens only beyond the aperture.  This identifies failure
of this chart.  Corollary~\ref{cor:alcubierre-maximality} and
Proposition~\ref{prop:flat-cap-rim} establish the invariant obstructions.

The finite-time wall sections in Figure~\ref{fig:reduced-patch} are
\(x=\pm\sqrt{1-y^2}\) and \(x=\pm\sqrt{1-\psi(y^2)}\) at \(z=0\),
using the profiles specified above.  The plotting script supplies the
smooth axial factors \(p=P(x)e^{\kappa t}\),
\(q=Q(x)e^{-\kappa t}\) used for the traveler in panel (a), including their
horizon limits.  The curve is a unit timelike geodesic on the axis \(\rho=0\),
with conserved Killing energy \(E\simeq1.1863>1\), so it belongs to the
core-accessible outgoing branch.  Its normalization and energy conservation
are checked across \(q=0\).


\subsection*{A qualification on earlier axial continuations}
For example, Eq.~(2.10) of Ref.~\cite{GonzalezDiaz2000}
reintroduces transverse terms after an axial reduction.  Off axis,
\(dr=(x\,dx+y\,dy+z\,dz)/r\); the axial identity \(dr=dx\) does not
establish a four-dimensional isometric pullback.
